\documentclass[UTF8]{ctexart}
\usepackage{xeCJK}
\usepackage{geometry}
\geometry{left=2cm,right=2cm,top=2.5cm,bottom=2.5cm}
\setlength{\headheight}{12.7pt}

\usepackage{titlesec}
\titleformat{\section}
  {\normalfont\large\bfseries}
  {\thesection}
  {1em}
  {}
\titlespacing*{\section}
  {0pt}
  {3.5ex plus 1ex minus .2ex}
  {2.3ex plus .2ex}

\usepackage{amsmath}
\usepackage{amssymb}
\usepackage{amsthm}
\usepackage{enumitem}
\setlist[enumerate]{itemsep=0pt,topsep=0pt,parsep=0pt,leftmargin=0pt,itemindent=2.5em}
\setlist[itemize]{itemsep=0pt,topsep=0pt,parsep=0pt,leftmargin=0pt,itemindent=2.5em}
\usepackage{fancyhdr}
\pagestyle{fancy}
\fancyhf{}
\usepackage{graphicx}
\usepackage{hyperref}
\usepackage{indentfirst}
\usepackage{lmodern}

\begin{document}
\setlength{\abovedisplayskip}{1pt}
\setlength{\belowdisplayskip}{1pt}

\section{Zorn 引理}

% 代数中神秘的Zorn引理.

\hypertarget{sec:定理1.1}{}
\noindent\textbf{定理1.1 (Zorn引理)}
\vspace{5pt}

任给\href{../01 序关系/02 偏序关系#nameddest=sec:定义1.1}{偏序结构}
$(A,\leqslant)$, 如果$A$的链(全序子集)都有上界, 那么$A$有极大元.

\vspace{7pt}\noindent{\kaishu 证明.}

假设$A$没有极大元, 即任意元素都有更大的元素.

对$A$的任意链$C$, 定义
\[
    m(C)=\{a\in A\mid a\text{\,\,是\,\,}C\text{\,\,的上界\,\,}
    \land a\notin C\,\}.\\[-3pt]
\]
显然$m(C)$非空: 取$C$的上界$b$, 再取比$b$更大的元素$a$, 就有$a\in m(C)$.

于是可以对$A$的任意链$C$选取一个$\varphi(C)\in m(C)$.

\vspace{7pt}
据此超限递归定义$\mathbf{On}$上的类映射$\mathbf{P}$ :
\vspace{3pt}
\[
    \forall\alpha\in\mathbf{On}:\quad\mathbf{P}(\alpha)=\left\{\begin{aligned}
        &\varphi\,(\,\mathbf{P}[\alpha]\,),\quad&&
        \mathbf{P}[\alpha]\text{\,\,是\,\,}A\text{\,\,的链\,},\\[-3pt]
        &A,\quad&&\mathbf{P}[\alpha]\text{\,\,不是\,\,}A\text{\,\,的链\,}.
    \end{aligned}\right.
\]

\vspace{7pt}\hspace{-1em}\noindent{\kaishu 命题1.}
任给序数$\alpha$, $\mathbf{P}[\alpha]$都是$A$的链.

对$\alpha\in\mathbf{On}$归纳证明.
任取序数$\alpha$, 假设对任意的$\beta<\alpha$, $\mathbf{P}[\beta]$是$A$的链, 此时
\[
    \mathbf{P}(\beta)=\varphi(\,\mathbf{P}[\beta]\,)\in m(\,\mathbf{P}[\beta]\,)
    \quad\implies\quad
    \mathbf{P}(\beta)\text{\,\,是\,\,}\mathbf{P}[\beta]\text{\,\,的上界\,\,}.
\]
然后任取$\gamma,\delta<\alpha$.

\begin{enumerate}
    \item 如果$\gamma<\delta$, 那么$\mathbf{P}(\gamma)\in\mathbf{P}[\delta]$,
    所以$\mathbf{P}(\gamma)\leqslant\mathbf{P}(\delta)$,
    \item 如果$\gamma=\delta$, 那么$\mathbf{P}(\gamma)=\mathbf{P}(\delta)$,
    \item 如果$\gamma>\delta$, 那么$\mathbf{P}(\delta)\in\mathbf{P}[\gamma]$,
    所以$\mathbf{P}(\gamma)\geqslant\mathbf{P}(\delta)$,
\end{enumerate}
所以$\mathbf{P}[\alpha]$是$A$的链.

\vspace{7pt}\hspace{-1em}\noindent{\kaishu 命题2.} $\mathbf{P}$是单射.

任给序数$\alpha,\beta$, 假设$\mathbf{P}(\alpha)=\mathbf{P}(\beta)$.
再假设$\alpha<\beta$, 则$\mathbf{P}(\alpha)\in\mathbf{P}[\beta]$, 同时由命题1得
\[
    \mathbf{P}(\beta)=\varphi(\,\mathbf{P}[\beta]\,)\in m(\,\mathbf{P}[\beta]\,)
    \quad\implies\quad\mathbf{P}(\beta)\notin\mathbf{P}[\beta],
\]
所以$\mathbf{P}(\alpha)\neq\mathbf{P}(\beta)$, 矛盾.

同理假设$\alpha>\beta$也矛盾. 所以$\alpha=\beta$.

\vspace{7pt}
由命题1, $\mathbf{P}(\alpha)$都属于$A$, 同时由命题2得$\mathbf{P}$是单射, 矛盾.
\hfill $\qedsymbol$





\newpage

\section{Hausdorff 极大原理}

\hypertarget{sec:定理2.1}{}
\noindent\textbf{定理2.1 (Hausdorff极大原理)}
\vspace{5pt}

偏序结构的链都可以扩张成极大链.

\vspace{7pt}\noindent{\kaishu 证明.}

任取偏序结构$(A,\leqslant)$和$A$的链$C$, 定义
\[
    \mathcal{S}=\{D\subset A\mid D\text{\,\,是\,\,}A\text{\,\,的链\,\,}
    \land D\supset C\},
\]
显然$\mathcal{S}$非空($C\in\mathcal{S}$).

考虑$\mathcal{S}$上的包含偏序$(\mathcal{S},\subset)$,
再任取$\mathcal{S}$的非空链$\mathcal{T}$.

\vspace{7pt}\hspace{-1em}\noindent{\kaishu 命题1.}
$\displaystyle\bigcup\mathcal{T}\subset A$,
$\displaystyle\bigcup\mathcal{T}\supset C$.

\vspace{3pt}
因为$\mathcal{T}\subset\mathcal{P}(A)$,
所以$\displaystyle\bigcup\mathcal{T}\subset A$.
再取$D\in\mathcal{T}$, 则有$\displaystyle C\subset D\subset\bigcup\mathcal{T}$.

\vspace{7pt}\hspace{-1em}\noindent{\kaishu 命题2.}
$\displaystyle\bigcup\mathcal{T}$是$A$的链.

\vspace{3pt}
对任意的$\displaystyle a,a'\in\bigcup\mathcal{T}$,
存在$D,D'\in\mathcal{T}$使得$a\in D,a'\in D'$. 因为$\mathcal{T}$是链, 所以:
\vspace{3pt}
\begin{enumerate}
    \item 如果$D\subset D'$, 那么$a,a'\in D'$, 再由$D'$是$A$的链可得
    $a\leqslant a'\lor a\geqslant a'$;
    \item 如果$D\supset D'$, 那么$a,a'\in D$, 同样由$D$是$A$的链可得
    $a\leqslant a'\lor a\geqslant a'$.
\end{enumerate}
\vspace{3pt}
综上可得$\displaystyle\bigcup\mathcal{T}$是$A$的链.

\vspace{7pt}
综上即得$\displaystyle\bigcup\mathcal{T}\in\mathcal{S}$,
显然$\displaystyle\bigcup\mathcal{T}$是$\mathcal{T}$的上界.

\vspace{7pt}
根据Zorn引理, $(\mathcal{S},\subset)$有极大元$M$, 于是:
\begin{enumerate}
    \item $M$是$A$的链且包含$C$,
    \item 如果$D$是$A$的链且包含$M$, 则$D\in\mathcal{S}$, 从而$D=M$.
    所以$M$是$A$的极大链. \hfill $\qedsymbol$
\end{enumerate}





\vspace{20pt}

\section{Tukey 引理}

\hypertarget{sec:定理3.1}{}
\noindent\textbf{定理3.1 (Tukey引理)}

任给非空集合$A$. 如果对任意的$x$, $x\in A$当且仅当
$x$的有限子集都属于$A$, 那么$(A,\subset)$有极大元.

\noindent{\kaishu 证明.}

任取$(A,\subset)$中的链$C$, 考虑$\displaystyle\bigcup C$.

对$\displaystyle\bigcup C$的任意非空有限子集$s$,
对任意的$x\in s$选取$f(x)\in C$使得$x\in f(x)$, 于是
\[
    \{f(x)\mid x\in s\}\subset C
\]
也是$(A,\subset)$的链, 因为它有限, 所以有最大元$m$. 从而对任意的$x\in s$有
\[
    f(x)\in C\quad\implies\quad f(x)\subset m\quad\implies\quad x\in m,
\]
即$s\subset m$, $s$是$m$的有限子集, 所以$s\in A$.

另外显然$\varnothing\in A$, 综上$\displaystyle\bigcup C$的有限子集都属于$A$,
所以$\displaystyle\bigcup C\in A$.

最后, 显然$\displaystyle\bigcup C$是$C$的上界.
于是$(A,\subset)$的链都有上界, 由Zorn引理得$(A,\subset)$有极大元. \hfill $\qedsymbol$

\end{document}
